% Artículo VII. Desarrollo aditivo, revisión 04, 10-09-2026.
% ARQUITECTURA_AUTORAL_PREEXISTENTE: transporte enriquecido, energía
% D_9^* W D_9 y compatibilidad por refinamiento, propietario S01,
% 02_dinamica_tres_hojas.tex, líneas 221--250.
% FORMALIZACION_NUEVA / CERTIFICADO_NUEVO: diferencial explícito de ese
% funcional, representante antihermítico, composición e inversión ponderada.
% Las variaciones pertenecen al espacio de realizaciones unitarias del
% transporte. No se presume que toda variación continua sea una ruta TPK.
% Su realización como corriente material de espín se trata por separado.
% \mathsf W_m es el peso energético; W de la sección anterior es soldadura.

\section{Respuesta variacional de la energía de memoria}
\label{vii:sec:variacion-memoria}

El transporte de una ruta identifica las fibras de sus extremos y conserva
la orientación y la memoria que intervienen en esa identificación. La energía
asociada compara el campo final con el campo inicial transportado. Su
diferencial cuantifica cómo responde esa comparación a una variación del
transporte. Esta sección obtiene la respuesta directamente del funcional
discreto y desarrolla su composición, covariancia e inversión.

\subsection{Funcional discreto y espacio de variaciones}
\label{vii:subsec:funcional-memoria}

En un nivel finito \(m\), sean \(\mathcal V_m\) los vértices de la
representación de rutas y \(\mathcal E_m^+\) una orientación de cada arista.
Cada vértice lleva una fibra hermítica finito-dimensional \(H_v\).
Definimos
\[
 \mathcal H_m=\bigoplus_{v\in\mathcal V_m}H_v,\qquad
 \mathcal K_m=\bigoplus_{a\in\mathcal E_m^+}H_{t(a)}.
\]
El producto hermítico es antilineal en la primera variable.
Para \(U_a:H_{s(a)}\to H_{t(a)}\) unitario, el operador de diferencias es
\begin{equation}
 (D_mf)_a=f_{t(a)}-U_af_{s(a)},\qquad
 E_m(f,U)=\langle D_mf,\mathsf W_mD_mf\rangle,\qquad
 \mathsf W_m=\mathsf W_m^*>0.
 \label{vii:eq:energia-memoria-discreta}
\end{equation}
Es la energía del término \(D_m^*\mathsf W_mD_m\) del generador discreto.
El peso \(\mathsf W_m\) puede acoplar distintas aristas. El operador \(W\)
de la soldadura de la sección~\ref{vii:sec:pantalla-respuesta} y este peso
energético tienen dominios y funciones diferentes.

El transporte \(U\) es una representación del estado enriquecido. El campo
\(f\) y el peso \(\mathsf W_m\) permanecen explícitos como argumentos
del funcional; cada aplicación material conserva las reglas que los
determinan. Se consideran colecciones diferenciables de las realizaciones
unitarias del transporte,
\[
 \dot U_a=A_aU_a,\qquad A_a^*=-A_a.
\]
Las variaciones físicamente utilizadas forman el subespacio que conserve
las condiciones de la realización correspondiente. El diferencial se
calcula primero en el espacio tangente unitario y después se restringe
a ese subespacio.

\begin{proposition}[Covector de respuesta al transporte]
\label{vii:prop:covector-memoria}
Escribamos
\[
 v_a=U_af_{s(a)},\qquad r=D_mf,\qquad q=\mathsf W_mr.
\]
A campo y peso fijos, el diferencial de la energía es
\begin{equation}
 \mathfrak j_m(A)=-2\operatorname{Re}
       \sum_{a\in\mathcal E_m^+}\langle q_a,A_av_a\rangle .
 \label{vii:eq:covector-memoria}
\end{equation}
Su representante para el producto real de Hilbert--Schmidt es
\begin{equation}
 \mathcal J_a=v_aq_a^*-q_av_a^*,\qquad
 \mathcal J_a^*=-\mathcal J_a,\qquad
 \mathfrak j_m(A)=
 \sum_a\operatorname{Re}\operatorname{Tr}(\mathcal J_a^*A_a).
 \label{vii:eq:representante-corriente-discreta}
\end{equation}
Para variaciones simultáneas de campo, transporte y peso se tiene
\begin{equation}
 \dot E_m=
 2\operatorname{Re}\sum_a
 \left\langle q_a,\dot f_{t(a)}-U_a\dot f_{s(a)}-A_av_a\right\rangle
 +\langle r,\dot{\mathsf W}_mr\rangle .
 \label{vii:eq:variacion-completa-memoria}
\end{equation}
\end{proposition}

\begin{proof}
La actualización de primer orden, escrita en el álgebra de números
duales con \(\epsilon^2=0\), da
\[
 r+\epsilon\dot r,\qquad
 \dot r_a=\dot f_{t(a)}-U_a\dot f_{s(a)}-A_av_a.
\]
El coeficiente de \(\epsilon\) en la energía es
\[
 \langle\dot r,\mathsf W_mr\rangle+
 \langle r,\mathsf W_m\dot r\rangle+
 \langle r,\dot{\mathsf W}_mr\rangle.
\]
La autoadjunción del peso produce
\eqref{vii:eq:variacion-completa-memoria}. Al fijar campo y peso queda
\eqref{vii:eq:covector-memoria}. Además,
\[
 \operatorname{Tr}(\mathcal J_a^*A_a)
 =v_a^*A_aq_a-q_a^*A_av_a.
\]
La antihermiticidad de \(A_a\) implica
\(v_a^*A_aq_a=-\overline{q_a^*A_av_a}\), de donde se obtiene
el representante anunciado.
\end{proof}

Así queda evaluada la respuesta: en cada arista se transporta el campo,
se calcula su diferencia con el campo final y se aplica el peso energético.
El producto antisimetrizado de esos dos vectores determina
\(\mathcal J_a\in\mathfrak u(H_{t(a)})\). Este operador representa,
mediante el producto real de Hilbert--Schmidt, el covector de respuesta
sobre las variaciones unitarias de la fibra de destino.

\subsection{Composición y transporte de la respuesta}
\label{vii:subsec:composicion-respuesta}

\begin{proposition}[Regla de composición sobre una ruta]
\label{vii:prop:composicion-respuesta}
Para una ruta formada por \(N\) transportes, sea
\(U_\gamma=U_N\cdots U_1\) y
\(B_k=U_N\cdots U_{k+1}\), con \(B_N=I\).
Si \(\dot U_k=A_kU_k\), entonces
\begin{equation}
 \dot U_\gamma U_\gamma^{-1}
 =\sum_{k=1}^N B_kA_kB_k^{-1}.
 \label{vii:eq:variacion-transporte-compuesto}
\end{equation}
Por tanto, un representante \(\mathcal J_\gamma\) en la fibra final
induce en el factor \(k\) la contribución
\(B_k^*\mathcal J_\gamma B_k\). Cuando un factor pertenece a varias
rutas del funcional, sus contribuciones se suman con las incidencias
que éste utiliza.
\end{proposition}

\begin{proof}
La regla del producto da
\[
 \dot U_\gamma=\sum_k
 U_N\cdots U_{k+1}A_kU_k\cdots U_1.
\]
Multiplicar por \(U_\gamma^{-1}\) elimina los factores a la derecha
de \(A_k\) y produce \eqref{vii:eq:variacion-transporte-compuesto}.
Por unitariedad y ciclicidad de la traza,
\[
 \operatorname{Re}\operatorname{Tr}
   (\mathcal J_\gamma^*B_kA_kB_k^{-1})
 =\operatorname{Re}\operatorname{Tr}
   ((B_k^*\mathcal J_\gamma B_k)^*A_k).
\]
La linealidad del diferencial demuestra la suma de contribuciones.
\end{proof}

Esta fórmula conserva la posición de cada operación dentro de la ruta.
El transporte posterior a una incidencia determina cómo se representa
su contribución en la fibra final; la respuesta local se recupera
transportándola a la fibra de esa incidencia.

\begin{proposition}[Covariancia bajo cambios unitarios de marco]
\label{vii:prop:covariancia-corriente-memoria}
Sean \(g_v\) cambios unitarios de marco, independientes del parámetro
de variación, y \(G_{\mathcal E}=\bigoplus_a g_{t(a)}\).
Bajo las transformaciones
\[
 f'_v=g_vf_v,\quad U'_a=g_{t(a)}U_ag_{s(a)}^{-1},\quad
 \mathsf W'_m=G_{\mathcal E}\mathsf W_mG_{\mathcal E}^{-1},\quad
 A'_a=g_{t(a)}A_ag_{t(a)}^{-1},
\]
se tiene
\begin{equation}
 \mathcal J'_a=g_{t(a)}\mathcal J_ag_{t(a)}^{-1},
 \qquad \mathfrak j'_m(A')=\mathfrak j_m(A).
 \label{vii:eq:covariancia-respuesta}
\end{equation}
\end{proposition}

\begin{proof}
Se obtiene \(r'=G_{\mathcal E}r\), \(q'=G_{\mathcal E}q\) y
\(v'_a=g_{t(a)}v_a\). La sustitución en
\eqref{vii:eq:representante-corriente-discreta} demuestra la primera
identidad. La invariancia del producto hermítico, o la ciclicidad de
la traza, demuestra la segunda.
\end{proof}

Cuando los cambios de marco dependan del parámetro, sus derivadas
actúan también sobre el campo y el peso. La expresión covariante
correspondiente es entonces el diferencial completo
\eqref{vii:eq:variacion-completa-memoria}.

\subsection{Inversión orientada y transporte del peso}
\label{vii:subsec:inversion-respuesta}

La inversión puede verse explícitamente en el caso de una energía
local por arista, \(\mathsf W_m=\bigoplus_a W_a\).
Para \(a:s\to t\), su inversa lleva
\[
 U_{\bar a}=U_a^{-1},\qquad
 r_{\bar a}=-U_a^{-1}r_a,\qquad
 W_{\bar a}=U_a^{-1}W_aU_a.
\]
Estas reglas conservan exactamente la energía de la arista.

\begin{proposition}[Compatibilidad diferencial con la inversión]
\label{vii:prop:inversion-respuesta}
Si \(W_a\) y los campos de ambos extremos permanecen fijos, entonces
\[
 A_{\bar a}=-U_a^{-1}A_aU_a,\qquad
 \dot W_{\bar a}=U_a^{-1}[W_a,A_a]U_a.
\]
La variación completa de la energía de la arista invertida coincide con
la original. En particular, si \([W_a,A_a]=0\), el covector parcial
de transporte satisface
\begin{equation}
 \mathfrak j_{\bar a}(U_a^{-1}A_aU_a)
 =-\mathfrak j_a(A_a).
 \label{vii:eq:imparidad-transporte}
\end{equation}
\end{proposition}

\begin{proof}
La derivada de \(U_a^{-1}\) es \(-U_a^{-1}A_a\), de donde resultan
las dos fórmulas. Escribiendo \(v_a=U_af_s\) y \(f_t=r_a+v_a\),
la parte debida exclusivamente al transporte invertido vale
\[
 \mathfrak j_{\bar a}(A_{\bar a})
 =-2\operatorname{Re}\langle r_a,W_aA_af_t\rangle
 =\mathfrak j_a(A_a)
   -2\operatorname{Re}\langle r_a,W_aA_ar_a\rangle.
\]
La variación del peso aporta
\[
 \langle r_{\bar a},\dot W_{\bar a}r_{\bar a}\rangle
 =\langle r_a,[W_a,A_a]r_a\rangle
 =2\operatorname{Re}\langle r_a,W_aA_ar_a\rangle.
\]
Ambos términos se compensan. Si el conmutador es nulo, esa contribución
desaparece; la linealidad en \(A_a\) da
\eqref{vii:eq:imparidad-transporte}.
\end{proof}

La energía se suma sobre una orientación por arista. Una suma sobre ambas
orientaciones lleva el factor \(1/2\). Esta convención conserva el mismo
funcional y, por tanto, la misma respuesta.

\subsection{Compatibilidad con el refinamiento}
\label{vii:subsec:refinamiento-respuesta}

\begin{theorem}[Transporte del diferencial por refinamiento]
\label{vii:thm:refinamiento-respuesta}
Sean \(J_m:\mathcal H_m\to\mathcal H_{m+1}\) y
\(K_m:\mathcal K_m\to\mathcal K_{m+1}\) aplicaciones fijas.
Para una colección \(C^1\), supongamos las identidades de refinamiento
\begin{equation}
 D_{m+1}(t)J_m=K_mD_m(t),\qquad
 K_m^*\mathsf W_{m+1}(t)K_m=\mathsf W_m(t).
 \label{vii:eq:familia-refinamiento-memoria}
\end{equation}
Entonces \(E_{m+1}(J_mf)=E_m(f)\).
Los covectores de transporte coinciden sobre las variaciones enlazadas
por la primera identidad. Cuando varían los pesos, sus contribuciones
al diferencial también coinciden.
\end{theorem}

\begin{proof}
Por sustitución,
\[
 \langle D_{m+1}J_mf,\mathsf W_{m+1}D_{m+1}J_mf\rangle
 =\langle D_mf,K_m^*\mathsf W_{m+1}K_mD_mf\rangle=E_m(f).
\]
Diferenciando la primera identidad de
\eqref{vii:eq:familia-refinamiento-memoria}, se obtiene
\(\dot D_{m+1}J_m=K_m\dot D_m\). Por consiguiente,
\[
 2\operatorname{Re}
 \langle D_{m+1}J_mf,\mathsf W_{m+1}\dot D_{m+1}J_mf\rangle
 =2\operatorname{Re}\langle D_mf,\mathsf W_m\dot D_mf\rangle.
\]
La derivada de la segunda identidad es
\(K_m^*\dot{\mathsf W}_{m+1}K_m=\dot{\mathsf W}_m\), que demuestra
la coincidencia de los términos de variación del peso.
\end{proof}

El enunciado utiliza la compatibilidad de la colección de transportes.
Si \(J_m\) depende del parámetro, la regla de la cadena conserva
también el término
\[
 2\operatorname{Re}
 \langle D_{m+1}J_mf,\mathsf W_{m+1}D_{m+1}\dot J_mf\rangle.
\]
La naturalidad obtenida transporta conjuntamente el funcional, sus
incidencias y su diferencial.

\subsection{Realización material de la respuesta}
\label{vii:subsec:realizacion-respuesta-discreta}

El resultado de esta sección es el covector explícito
\(\mathfrak j_m\), con su representante \(\mathcal J_a\).
La realización material compone ese covector con la variación del
transporte inducida por la conexión y con la normalización de la acción.
Si \(\mathcal L_m\) denota ese mapa diferencial, el covector sobre
variaciones de conexión es \(\mathcal L_m^*\mathfrak j_m\).
La regla de composición
\eqref{vii:eq:variacion-transporte-compuesto} evalúa su parte
correspondiente a productos finitos de transportes.

En la sección~\ref{vii:sec:corriente-respuesta-cartan}, la corriente
\(\sigma_{IJ}\) es una tres-forma definida por la variación de la acción
material. La identificación entre ambas realizaciones conserva el campo,
la representación, el peso, la normalización dimensional y el paso al
límite utilizados por esa acción. Estos datos permiten evaluar la
contracción que interviene en la densidad torsional. El desarrollo
presente aporta el diferencial discreto que se compone en ese cálculo.
